\documentclass[11pt,a4paper]{article}
\usepackage[utf8]{inputenc}
\usepackage{amsmath,amssymb}
\usepackage{geometry}
\geometry{margin=2.5cm}
\usepackage{listings}
\usepackage{xcolor}
\usepackage{hyperref}

\lstset{
  language=C++,
  basicstyle=\ttfamily\small,
  keywordstyle=\color{blue}\bfseries,
  commentstyle=\color{green!60!black},
  stringstyle=\color{red!70!black},
  numbers=left,
  numberstyle=\tiny\color{gray},
  numbersep=8pt,
  frame=single,
  breaklines=true,
  tabsize=4,
  showstringspaces=false
}

\title{IOI 1994 -- The Clock}
\author{Solution}
\date{}

\begin{document}
\maketitle

\section{Problem Statement}

There are 9 clocks arranged in a $3 \times 3$ grid. Each clock points to 12, 3, 6, or 9
(represented as 0, 1, 2, 3 respectively, where 0 means 12 o'clock).
There are 9 possible moves, each of which rotates a specific subset of clocks
by 90 degrees clockwise. The goal is to make all clocks point to 12 (all zeros)
using the minimum number of moves.

\medskip
\noindent The 9 moves affect the following clocks (numbered 1--9 in row-major order):
\begin{itemize}
  \item Move 1: clocks 1, 2, 4, 5
  \item Move 2: clocks 1, 2, 3
  \item Move 3: clocks 2, 3, 5, 6
  \item Move 4: clocks 1, 4, 7
  \item Move 5: clocks 2, 4, 5, 6, 8
  \item Move 6: clocks 3, 6, 9
  \item Move 7: clocks 4, 5, 7, 8
  \item Move 8: clocks 7, 8, 9
  \item Move 9: clocks 5, 6, 8, 9
\end{itemize}

\section{Solution Approach}

\subsection{Key Observation}
Since each clock has only 4 states (mod 4), and each move applied 4 times returns
all affected clocks to their original positions, each move needs to be applied
at most 3 times. This means the total state space is $4^9 = 262{,}144$ states,
which is small enough for BFS.

\subsection{Algorithm: BFS}
\begin{enumerate}
  \item Encode the 9 clock values as a single state (e.g., treating each clock as
        a base-4 digit).
  \item Starting from the initial state, perform BFS.
  \item Each state has 9 neighbors (one per move).
  \item The target state is $0$ (all clocks at 12).
  \item BFS guarantees the shortest sequence of moves.
\end{enumerate}

\subsection{Alternative: Brute Force with 4 Loops}
Since each of the 9 moves is applied 0--3 times, we can enumerate all $4^9$ combinations.
This is feasible and simpler to implement.

\section{C++ Solution}

\begin{lstlisting}
#include <cstdio>
#include <cstring>
#include <queue>
#include <map>
using namespace std;

// Clocks affected by each move (0-indexed)
int moves[9][5] = {
    {0, 1, 3, 4, -1},   // Move 1: clocks 1,2,4,5
    {0, 1, 2, -1, -1},  // Move 2: clocks 1,2,3
    {1, 2, 4, 5, -1},   // Move 3: clocks 2,3,5,6
    {0, 3, 6, -1, -1},  // Move 4: clocks 1,4,7
    {1, 3, 4, 5, 7},    // Move 5: clocks 2,4,5,6,8
    {2, 5, 8, -1, -1},  // Move 6: clocks 3,6,9
    {3, 4, 6, 7, -1},   // Move 7: clocks 4,5,7,8
    {6, 7, 8, -1, -1},  // Move 8: clocks 7,8,9
    {4, 5, 7, 8, -1}    // Move 9: clocks 5,6,8,9
};

int clk[9];

int encode(int c[]) {
    int s = 0;
    for (int i = 0; i < 9; i++)
        s = s * 4 + c[i];
    return s;
}

void decode(int s, int c[]) {
    for (int i = 8; i >= 0; i--) {
        c[i] = s % 4;
        s /= 4;
    }
}

void applyMove(int c[], int m) {
    for (int k = 0; k < 5 && moves[m][k] != -1; k++)
        c[moves[m][k]] = (c[moves[m][k]] + 1) % 4;
}

int main() {
    for (int i = 0; i < 9; i++) {
        int x;
        scanf("%d", &x);
        // Convert: 12->0, 3->1, 6->2, 9->3
        clk[i] = (x / 3) % 4;
    }

    // BFS
    int start = encode(clk);
    int target = 0; // all zeros

    if (start == target) {
        printf("0\n");
        return 0;
    }

    map<int, int> dist;
    map<int, pair<int, int>> parent; // state -> (prev_state, move)
    queue<int> q;
    q.push(start);
    dist[start] = 0;

    while (!q.empty()) {
        int u = q.front(); q.pop();
        int d = dist[u];

        for (int m = 0; m < 9; m++) {
            int tmp[9];
            decode(u, tmp);
            applyMove(tmp, m);
            int v = encode(tmp);
            if (dist.find(v) == dist.end()) {
                dist[v] = d + 1;
                parent[v] = {u, m};
                if (v == target) {
                    // Reconstruct path
                    int path[1000], plen = 0;
                    int cur = target;
                    while (cur != start) {
                        path[plen++] = parent[cur].second + 1;
                        cur = parent[cur].first;
                    }
                    for (int i = plen - 1; i >= 0; i--) {
                        if (i < plen - 1) printf(" ");
                        printf("%d", path[i]);
                    }
                    printf("\n");
                    return 0;
                }
                q.push(v);
            }
        }
    }

    return 0;
}
\end{lstlisting}

\section{Complexity Analysis}

\begin{itemize}
  \item \textbf{Time complexity:} $O(4^9 \times 9) = O(2{,}359{,}296)$.
        There are at most $4^9 = 262{,}144$ distinct states, and from each state
        we try 9 moves. This is very fast in practice.
  \item \textbf{Space complexity:} $O(4^9)$ for the BFS visited set and parent map.
\end{itemize}

\end{document}
